\subsection{Exact checks of the cost of discarded reports}
\label{sec:summaryexperiment}

The finite-observation planning gap and the loss of an omitted observation
are different quantities. This experiment checks
\cref{thm:summarysandwich,prop:binarysummary} on a finite tree where both can
be evaluated exactly. It also shows why a fine belief grid cannot recover
information that the controller never receives.

\subsubsection*{Controlled liquidity with a coarser controller}

Use the seven instructions, transitions, fills and charges in
\eqref{eq:partialexperimentkernel}--\eqref{eq:partialexperimentcost}, with
three decisions, two quantity units and $b_0=1/2$. The richer agent observes
both the fill and public signal. The implemented summary controller retains
the quantity and a rounded filter driven only by fills; it discards the
public signal. Its update uses the kernel obtained by summing
\eqref{eq:partialexperimentkernel} over $y$. It retains the same exact
quantity mask as the richer agent.

All positive-probability fine-report histories under all feasible actions
are enumerated with rational arithmetic. Histories are merged only when
their stage, quantity, controller memory and exact full posterior agree.
At each summary, the smallest and largest of these full posteriors give
\eqref{eq:posteriorfiber}. The endpoint recursions then produce $L,U$ and
the upper-recursion minimizing policy $\sigma$. Independently, full-posterior
dynamic programming computes $V^{\rm rich}$, and fill-history dynamic
programming computes the exact coarse-history optimum $V^{\rm coarse}$.
The latter retains the entire fill history and has no rounding error.
A separate forward calculation propagates actual hidden-state probabilities
under the fixed $\sigma$ to obtain its physical cost $J_\sigma$.

\begin{table}[tb]
\centering\small
\caption{Exact finite-tree observation-reduction checks, in USD.
$z_0/z_1$ denotes the two probabilities of public signal one given next
hidden state; $m$ is the number of uniform belief subintervals.
The two exact optima and $L_0$ are displayed downward to six decimals;
$J_\sigma$ and $J_\sigma-L_0$ are displayed upward. The machine-readable
results retain every exact rational value.}
\label{tab:summarycertificate}
\begin{tabular}{rrrrrrr}
\toprule
$z_0/z_1$ & $m$ & $V^{\rm rich}$ & $V^{\rm coarse}$ &
$J_\sigma$ & $L_0$ & Certified gap\\
\midrule
\inputtable{tables/observation_summary_rows.tex}
\bottomrule
\end{tabular}
\end{table}

The first row uses the original public-signal probabilities. Its richer
optimum is the previously checked
$232172138011/156250000000=1.4859016832704$.
Discarding the signal raises the optimal expected cost by exactly
$1293658149/156250000000=0.0082794121536$ USD, before any rounding or
summary-policy restriction. The remaining rows strengthen only the report,
to probabilities $0.05$ and $0.95$. The coarse model is unchanged because
it integrates out that report. Its information loss rises to exactly
\[
 V^{\rm coarse}-V^{\rm rich}
 =\frac{205859933299369}{2500000000000000}
 =0.0823439733197476\ \mathrm{USD}.
\]
Increasing $m$ from eight to sixteen slightly improves the
certificate while leaving the selected controller's initial cost unchanged.

These are valid but conservative certificates: the stagewise interval
extrema can combine posterior histories that one physical trajectory cannot
jointly realize. The upper-recursion policy need not be optimal among
summary policies. The reported gap uses the sharper physical evaluation
$J_\sigma-L_0$, not $U_0-L_0$, and is compared with the exact achieved gap
$J_\sigma-V^{\rm rich}$ in the audit output.

There are $1,24,324,3456$ distinct full-posterior/memory states at the four
stages in each case. Every one of the 349 nonterminal states satisfies
$L\le V^{\rm rich}\le J^\sigma\le U$ in exact arithmetic. The separate
forward hidden-state evaluation equals the history-based evaluation exactly
and conserves probability at every stage. The maximum discrepancy between
the full posterior and rounded memory exceeds $0.43$ even in the first
case, whereas the one-step rounding radius is $1/16$. This directly checks
that those two errors cannot be identified.

\subsubsection*{A transparent probe example and the sufficiency limit}

For a separate two-decision example, one unit remains and good/bad liquidity
is static with equal prior probabilities. At the first decision the trader
may wait at zero cost or pay $0.02$ USD for a report that identifies the
state correctly with probability $0.9$. At the second decision, a market
instruction completes at cost $1$ USD. Alternatively, a passive instruction
fills at zero cost with probability $0.9$ in good liquidity and $0.1$ in bad
liquidity; any failure is completed at cost $2$ USD. Every instruction is
feasible and all residuals are completed.

Writing $b$ for the probability of good liquidity in this example, the
passive expected cost is $1.8-1.6b$. A retained favorable report gives
$b=0.9$ and passive cost $0.36$; an unfavorable report gives $b=0.1$ and
selects the market instruction. Thus the richer agent probes and has value
\[
 V^{\rm rich}=0.02+\tfrac12(0.36+1)=0.70.
\]
An agent discarding the report has exact posterior $b=1/2$, optimally waits,
and has value $V^{\rm coarse}=1$. Its information loss is $0.30$ USD even
with an exact coarse filter and no quantization. If the summary remembers
whether it probed but discards the report, its post-probe posterior envelope
is $[0.1,0.9]$. The endpoint recursions give
\[
 L_0=\min\{1,0.02+0.36\}=0.38,\qquad
 U_0=\min\{1,0.02+1\}=1.
\]
The resulting $0.62$ USD certificate covers the $0.30$ USD actual loss.
Retaining the report makes every conditional envelope a singleton; both
recursions then equal $0.70$, as required by
\cref{cor:summarysufficiency}. These values are checked as exact fractions.

\subsubsection*{Action certificates and decision sufficiency}

The same three finite trees also check
\cref{thm:summaryactioncertificate}. Every one of the 1,237 feasible
state--action pairs per case satisfies both action bounds in
\eqref{eq:summaryactionbounds} and the Bellman-advantage bound
\eqref{eq:summaryactionexcess}, using exact fractions. The resulting regret
recursion bounds the selected controller at all 349 nonterminal states.
Across the three cases, 648, 409 and 417 strict action-screening comparisons
identify only truly suboptimal actions. These counts include the different
full histories in a common summary fiber.

\begin{table}[tb]
\centering\small
\caption{Two certificates for the same implemented controller in the
three controlled-liquidity cases, in USD. The cost certificate is
$J_\sigma-L_0$ using physical evaluation; the action certificate is
$R_0^\sigma$. The combined bound takes their minimum. All displayed gaps
and bounds are rounded upward to six decimals from exact rational values.}
\label{tab:actioncertificate}
\begin{tabular}{rrrrrr}
\toprule
$z_0/z_1$ & $m$ & Actual gap & Cost certificate & Action certificate & Combined\\
\midrule
\inputtable{tables/action_separation_rows.tex}
\bottomrule
\end{tabular}
\end{table}

Here the cost certificate is sharper in every row of
\cref{tab:actioncertificate}; the action recursion is not a uniformly
tighter replacement. A separate two-stage case shows the distinction.
At the first stage the only feasible action costs $0.02$ USD and produces
a perfect report of a static, equally likely binary state. The richer agent
retains this report and the summary controller discards it. At the second
stage both terminal instructions complete the remaining unit by assumption.
Instruction $L$ has all-in charges $0.10$ and $0.20$ USD in the two states;
instruction $M$ costs $0.50$ USD in either state. These supplied charges
define a finite execution example, not an estimated exchange fill law.

The summary's posterior envelope after the report is $[0,1]$. Its action
bounds are $Q^-(L)=0.10$, $Q^+(L)=0.20$ and
$Q^-(M)=Q^+(M)=0.50$. Hence $e(L)=0$ and $e(M)=0.40$.
The fixed controller choosing $L$ has
\[
 L_0=0.12,\qquad J_L=V_0^{\rm rich}=0.17,\qquad U_0=0.22.
\]
The physically evaluated cost certificate is $0.05$ USD, but the action
certificate is exactly zero. It establishes optimality despite the
non-singleton envelope and nonzero cost interval. The richer conditional
value at the final decision is $0.10$ or $0.20$, so the summary is
\emph{decision sufficient} here without determining that conditional value.

For a randomized controller choosing $M$ with probability $1/4$, physical
cost is $0.2575$ USD and true regret is $0.0875$ USD. The cost certificate
is $0.1375$ USD while the action certificate is $0.10$ USD. The script
checks these quantities and the conditional bounds in both hidden states
as exact fractions. This also checks the action average in
\eqref{eq:summaryactionrecursion}; randomization is evaluated as implemented.

The script \texttt{scripts/observation\_summary\_checks.py} generates both
tables and stores the exact fractions, state counts and conditional checks.
The same script checks the Gaussian sign symmetry in
\cref{thm:ledgerlower} and the nonuniform cell created by inserting the
initial prior $1/3$ into an eight-subinterval grid. These are finite
verification cases supporting the written proofs; they do not constitute a
machine-checked proof of the general conditional-envelope theorem.
